\chapter{Deneysel Kurulum}
\label{ch:deneysel_kurulum}

Bu bölümde tezin deneysel değerlendirmesinde kullanılan veri kümeleri, donanım/yazılım altyapısı ve eğitim parametreleri açıklanmaktadır.

\section{Veri Kümeleri}
\subsection{NIH Chest X-Ray 14}
Çalışmada pnömotoraks patolojisinin tespiti için NIH Chest X-Ray 14 veri seti kullanılmıştır. Veri setinden elde edilen röntgen görüntüleri CLAHE (Contrast Limited Adaptive Histogram Equalization) ve normalizasyon ön işlemlerinden geçirilmiştir.

\subsection{CheXpert}
Geliştirilen kademeli sistemin genelleştirilebilirliğini test etmek ve sıfır-atışlı (zero-shot) çapraz veri kümesi performansını ölçmek amacıyla Stanford CheXpert veri kümesi kullanılmıştır.

\section{Donanım ve Yazılım Altyapısı}
Eğitim ve test işlemleri Apple M3 Max (MPS) ve NVIDIA RTX 4090 GPU donanımları üzerinde gerçekleştirilmiştir. Yazılım altyapısı olarak PyTorch (v2.3.0), FastAPI (v0.111.0), ve React (v18.3.1) kullanılmıştır. Deneylerin takibi MLflow üzerinden gerçekleştirilmiştir.

\section{Eğitim Parametreleri}
Tier 1 ve Tier 2 modellerinin eğitiminde kullanılan hiper-parametreler Tablo \ref{tab:hiperparametreler} içinde özetlenmiştir.

\begin{table}[htbp]
\centering
\caption{Model Eğitim Hiper-parametreleri}
\label{tab:hiperparametreler}
\begin{tabular}{lcc}
\toprule
\textbf{Parametre} & \textbf{Tier 1 (MobileNetV2)} & \textbf{Tier 2 (EfficientNetB4)} \\
\midrule
Öğrenme Oranı (Learning Rate) & $10^{-4}$ & $5 \times 10^{-5}$ \\
Batch Boyutu & $64$ & $32$ \\
Optimizasyon Algoritması & Adam & AdamW \\
Epoch Sayısı & $50$ & $30$ \\
Görüntü Boyutu & $224 \times 224$ & $380 \times 380$ \\
\bottomrule
\end{tabular}
\end{table}
